\documentclass[11pt]{article}

\usepackage[T1]{fontenc}
\usepackage{lmodern}
\usepackage{amsmath,amssymb,amsthm}
\usepackage[margin=1in]{geometry}
\usepackage{microtype}
\usepackage[hidelinks,pdftitle={The Exact Ternary Value ASP(6,4,3)=12},pdfauthor={Xiaojian Zeng}]{hyperref}
\usepackage{booktabs}
\usepackage{enumitem}

\newtheorem{theorem}{Theorem}
\newtheorem{lemma}[theorem]{Lemma}
\newtheorem{proposition}[theorem]{Proposition}
\newcommand{\F}{\mathbb F}
\newcommand{\ASP}{\operatorname{ASP}}
\newcommand{\ASB}{\operatorname{ASB}}

\title{The Exact Ternary Value $\ASP(6,4,3)=12$}
\author{Xiaojian Zeng\\
  \small PBC School of Finance, Tsinghua University, 100190\\
  \small \texttt{xiaojianzeng@foxmail.com}}
\date{}

\begin{document}
\maketitle

\begin{abstract}
For a rank-$k$ matrix over $\F_q$, the $t$-all-symbol PIR property asks
for $t$ pairwise disjoint recovery sets for every column.  The closest
available bounds give
$11\leq\ASP(6,4,3)\leq\ASB(6,4,3)\leq12$.
We prove that the lower endpoint is impossible and hence that
$\ASP(6,4,3)=12$.  The proof reduces every putative ternary
$6\times11$ matrix, without loss, to a normalized systematic matrix and
encodes the remaining finite problem as a $7{,}428$-variable,
$34{,}731$-clause CNF.  We prove both directions of the encoding and
certify unsatisfiability by a replayed DRUP derivation.  A separate exact
kernel enumeration tests the semantics of the encoding on $25$ fixed
matrices and all their targets; it is an audit of the encoding, not a
substitute for the refutation.
\end{abstract}

\section{Statement and context}

Coding can replace replication in distributed private information
retrieval, leading to PIR codes with disjoint recovery groups
\cite{FazeliVardyYaakobi2015}.  Boruchovsky, Gruica, Niemann, and Yaakobi
introduced the all-symbol parameters used here
\cite{BoruchovskyEtAl2026}.  If
$G=[g_1,\ldots,g_n]\in\F_q^{k\times n}$ has rank $k$, a set
$R\subseteq[n]$ \emph{recovers} $g_i$ when
\[
  g_i\in\operatorname{span}_{\F_q}\{g_j:j\in R\}.
\]
The matrix is $t$-all-symbol PIR when every column has $t$ pairwise
disjoint recovery sets, and $\ASP(k,t,q)$ is the minimum possible
length.  The batch parameter $\ASB(k,t,q)$ imposes the stronger
simultaneous recovery requirement.

Theorem~23 of \cite{BoruchovskyEtAl2026}, specialized to $k=6$ and
$q=3$, gives
\begin{equation}\label{eq:published-interval}
  11\leq\ASP(6,4,3)\leq\ASB(6,4,3)\leq12.
\end{equation}
That paper reports exhaustive nonexistence for the corresponding binary
length-$11$ cell, not for the ternary cell.  Thus the precise remaining
task in \eqref{eq:published-interval} is to decide whether a rank-six
$6\times11$ matrix over $\F_3$ with the four-all-symbol PIR property
exists.  We prove that it does not.  Our originality statement is
deliberately limited: a bounded search located no later primary
resolution, but this does not establish priority.

\begin{theorem}\label{thm:main}
There is no rank-six matrix in $\F_3^{6\times11}$ having four pairwise
disjoint recovery sets for every column.  Consequently,
\[
  \ASP(6,4,3)=12.
\]
\end{theorem}

\section{Lossless normalization}

We first reduce the quantified matrix problem to the exact domain used
by the certificate.

\begin{lemma}[Three nontrivial recoveries]\label{lem:three}
Let $g_i\neq0$.  The column $g_i$ has four pairwise disjoint recovery
sets if and only if it has three pairwise disjoint recovery sets that
avoid $i$, in addition to the singleton $\{i\}$.
\end{lemma}

\begin{proof}
The reverse implication is immediate.  Conversely, among four disjoint
sets at most one contains $i$.  If one does, replace it by $\{i\}$ and
retain the other three.  If none does, retain any three and add
$\{i\}$.
\end{proof}

\begin{lemma}[Systematic normalized form]\label{lem:normalization}
If a matrix forbidden by Theorem~\ref{thm:main} exists, then one exists
in the form
\[
  G=(I_6\mid A),\qquad A\in\F_3^{6\times5},
\]
with the following properties.
\begin{enumerate}
  \item[(i)] Every column of $A$ is nonzero and its first nonzero entry is
        $1$.
  \item[(ii)] The first entries of the first three columns of $A$ are $1$.
  \item[(iii)] Every column of $A$ whose first entry is $1$ precedes every
        column whose first entry is $0$.
\end{enumerate}
\end{lemma}

\begin{proof}
No column of the original matrix is zero.  Indeed, deleting a zero
column preserves rank, and deleting that coordinate from all recovery
sets preserves their spans and disjointness.  This would produce a
length-$10$ example, contradicting the published lower bound in
\eqref{eq:published-interval}.

Choose six basis columns, permute them first, and left-multiply by the
inverse basis matrix.  Every recovery equation is preserved, so the
matrix becomes $(I_6\mid A)$.  Each remaining column is nonzero.
Multiplying any individual column $g_j$ by a scalar $d_j\neq0$ also
preserves the property: the equation
$g_i=\sum_j\alpha_jg_j$ becomes
\[
 d_i g_i=\sum_j \alpha_j\frac{d_i}{d_j}(d_jg_j).
\]
We may therefore scale each column of $A$ so that its first nonzero
entry is $1$.

Apply Lemma~\ref{lem:three} to the systematic target $e_1$.  Each of
its three disjoint target-avoiding recoveries must use a column of $A$
with nonzero first entry, because $e_2,\ldots,e_6$ have zero first
entry.  Disjointness makes these three parity columns distinct.  Their
normalized first entries are $1$.  A permutation of the five columns
of $A$ now gives (ii), and a further permutation within the two
first-entry classes gives (iii).  All operations used are invertible or
are column permutations, so no candidate is lost.
\end{proof}

\section{Exact CNF equivalence}

We describe the formula generated by the ancillary program
\texttt{systematic\_sat\_search.py}.  Every field variable is represented
by three Boolean literals, one for each value in $\{0,1,2\}$, with one
at-least-one and three pairwise at-most-one clauses.  Thus every such
trit has exactly one value.

The $30$ trits $a_{r,p}$ encode $A$.  Clauses require each parity
column to be nonzero, make its first nonzero entry equal to $1$, set
$a_{1,1}=a_{1,2}=a_{1,3}=1$, and put all first-row $1$ entries before
all first-row $0$ entries.  These are precisely the conventions in
Lemma~\ref{lem:normalization}.

For every target $t\in[11]$ and recovery label $s\in[3]$, trits
$c_{t,s,j}$ encode a coefficient on column $j$.  The formula imposes
\begin{align}
  c_{t,s,t}&=0,\label{eq:avoid}\\
  \operatorname{supp}(c_{t,s,\cdot})\cap
  \operatorname{supp}(c_{t,s',\cdot})&=\varnothing
    \quad(s\neq s'),\label{eq:disjoint}\\
  \sum_{j=1}^{6}c_{t,s,j}e_j+
  \sum_{p=1}^{5}c_{t,s,6+p}a_{\cdot,p}&=g_t
    \quad\text{in }\F_3^6.\label{eq:recovery}
\end{align}
For each support an exact-one family of ``minimum'' variables chooses
its least nonzero coordinate.  The three minima are required to be
strictly increasing.  Hence every support is nonempty and the constraint
only orders the three otherwise interchangeable recoveries.

Equation~\eqref{eq:recovery} is expanded row by row.  Multiplication and
addition gates contain an implication for each of the nine possible
input pairs in $\F_3^2$.  Together with exact-one input and output trits,
these exhaustive truth tables make every gate functional and exact;
they are not a one-way relaxation.  The resulting formula $F$ has
$7{,}428$ variables and $34{,}731$ clauses.

\begin{proposition}[Encoding equivalence]\label{prop:equivalence}
The CNF $F$ is satisfiable if and only if a rank-six
$6\times11$ matrix over $\F_3$ has the four-all-symbol PIR property.
\end{proposition}

\begin{proof}
From a satisfying assignment, decode the exact-one trits.  The matrix
is $(I_6\mid A)$ and therefore has rank six; its parity columns satisfy
the normalization clauses.  For each target, equations
\eqref{eq:avoid}--\eqref{eq:recovery} give three nonempty, pairwise
disjoint, target-avoiding coefficient supports that recover the target.
All target columns are nonzero.  Lemma~\ref{lem:three} supplies a fourth
recovery, namely the singleton target.

Conversely, normalize any putative matrix by
Lemma~\ref{lem:normalization}.  By Lemma~\ref{lem:three}, choose three
disjoint target-avoiding recoveries for each target and choose
coefficients expressing the target in each span.  Delete coordinates
whose coefficient is zero.  The resulting supports remain disjoint and
are nonempty because the targets are nonzero.  Order the three supports
by their distinct least coordinates, assign all coefficient trits, and
assign each arithmetic gate its actual value in $\F_3$.  Every matrix,
minimum, disjointness, and truth-table clause is then satisfied.
\end{proof}

\section{The checked refutation}

The clausal-certificate conventions used below are standard; see
\cite{WetzlerHeuleHunt2014,Heule2016}.  A clause $C$ has the reverse unit
propagation (RUP) property with respect to a clause database $D$ when
unit propagation on
$D\cup\{\{\lnot\ell\}:\ell\in C\}$ reaches a conflict.  In that case
the clause $C$ is entailed by $D$.

The supplied file \texttt{systematic\_search\_glucose.drup} was produced
by Glucose4, from the solver family introduced in
\cite{AudemardSimon2009}.  It has $833{,}442$ lines: $426{,}775$ clause
additions and $406{,}667$ deletions.  The short checker
\texttt{replay\_drup.py} parses the fixed CNF and proof, checks every
addition by RUP, and reaches the empty clause on the final line.  It
uses Minisat22 only as a unit-propagation engine and performs no
satisfiability search.

For completeness, ignoring the deletion instructions is sound here.
Inductively, every retained added clause is entailed by the input CNF.
The retained database is therefore equivalent to the input with respect
to satisfiability.  If unit propagation refutes the negation of the next
clause, that clause too is entailed by the input.  The accepted empty
clause consequently proves that $F$ is unsatisfiable.  This argument
establishes the following certificate proposition.

\begin{proposition}[Certified finite fact]\label{prop:unsat}
The formula $F$ in Proposition~\ref{prop:equivalence} is unsatisfiable.
\end{proposition}

The release manifest binds the exact CNF, trace, checker, and replay
report.  From the project root, the decisive check is reproduced by
\begin{verbatim}
python3 evidence/replay_drup.py evidence/systematic_search.cnf \
  evidence/systematic_search_glucose.drup --solver minisat22
\end{verbatim}
The output must report acceptance, all $426{,}775$ RUP additions, and an
empty clause on line $833{,}442$.

Combining Propositions~\ref{prop:equivalence} and \ref{prop:unsat}
excludes the length-$11$ matrix.  Equation~\eqref{eq:published-interval}
then proves Theorem~\ref{thm:main}.

\section{Independent semantic tests and limits}

The separate program \texttt{validate\_sat\_encoding\_broad.py} checks
the meaning of fixed-matrix instances without reusing the CNF arithmetic
as its semantic oracle.  For $G=(I_6\mid A)$, every kernel vector is
uniquely
\[
  (-Ax,x),\qquad x\in\F_3^5.
\]
A target-avoiding recovery of column $t$ is equivalent, after scaling,
to a projective kernel direction whose $t$th coordinate is nonzero:
if $Mc=g_t$ and $c_t=0$, then $c-e_t\in\ker M$ has $t$th coordinate
$-1$, and the converse follows by reversing this construction.
The oracle enumerates all
$(3^5-1)/(3-1)=121$ projective directions, constructs their exact
supports, and finds the maximum number of disjoint supports by a
bit-mask dynamic program.

This oracle was compared with the targetwise CNF on $25$ normalized
matrices, all $11$ targets, and two SAT backends, for $550$ status checks
in total.  All statuses agreed, with both satisfiable and unsatisfiable
target instances represented.  These tests could expose a semantic
encoding error but do not range over every $A$ and do not prove global
unsatisfiability; Proposition~\ref{prop:unsat} rests on the DRUP
certificate.  The separation between discovery and an exact verifier
that reconstructs the finite predicates also appears in certified
finite ternary searches of Zeng, Liu, Ratnavelu, and Huat
\cite{ZengEtAl2026TernaryCertificates}.

The result concerns exactly $(k,t,q,n)=(6,4,3,11)$.  It does not classify
all extremal length-$12$ matrices, and it makes no assertion for other
fields or dimensions.  A second model context audited the frozen
reduction and replayed the certificate; this is a useful workflow check,
not external peer review.

\section*{Computational and AI assistance}

OpenAI Codex assisted with the bounded literature search, the design and
implementation of the normalization-aware SAT encoding, the exact
semantic validator, certificate replay and audit organization, and
initial manuscript preparation.  Language-model output is not used as
mathematical evidence.  The theorem-bearing evidence consists of the
explicit reduction above, the fixed CNF and clausal derivation, and the
replayable exact checks.  No proof assistant was used.

\bibliographystyle{plain}
\bibliography{references}

\end{document}
